\begin{frame}{ROC-AUC vs PR-AUC: 언제 무엇을 보아야 하는가}
  \begin{columns}[T]
    \begin{column}{0.5\textwidth}
      \kb{ROC-AUC}
      \begin{itemize}
        \item TPR-FPR 쌍을 임계값별로 그리는 곡선.
        \item \textbf{클래스 균형이 거의 맞을 때} 좋은 스크리닝
          지표.
        \item 무작위 0.5, 완벽 1.0.
        \item \textbf{불균형이 무너지면} ``무조건 0'' 모델의
          ROC-AUC가 여전히 0.5라, 0.15\% 양성에서 ROC-AUC=0.51도
          ``무작위보다 낫다''로 읽혀 \textbf{과대평가}되기 쉬움.
      \end{itemize}
    \end{column}
    \begin{column}{0.5\textwidth}
      \kb{PR-AUC}
      \begin{itemize}
        \item \textbf{클래스 불균형이 심할 때} 더 ``정직한'' 지표 ---
          기준선이 양성 비율이라, 양성 15\%에서 PR-AUC=0.151이면
          ``무조건 0과 거의 같은 모델''이 바로 읽힘.
        \item ROC-AUC=0.51은 ``무작위보다 약간 낫다''로 읽혀
          \textbf{절대 척도가 다름}.
      \end{itemize}
    \end{column}
  \end{columns}
  \vskip0.6em
  \begin{block}{실무 권장 흐름}
    \kb{ROC-AUC로 후보 모델 스�리닝 $\to$ PR-AUC와 비용 곡선으로
    최종 임계값 선택}.
    \\ ROC-AUC는 ``순서 짓기 능력''을 빠르게 비교하고, PR-AUC +
    비용 최소화는 ``이 문제를 실제로 풀 때 어떤 임계값이 가장
    이익''을 결정.
  \end{block}
\end{frame}
